\chapter{渐近展开.}

各阶“无穷小”（或“无穷大”）之间的区分，从关于微分学的最早著作、例如费马的著作起，就已经隐含地出现；它随着牛顿和莱布尼茨的“高阶差分”理论而变成完全自觉的；而且人们随即便观察到，在最简单的情形，形如 \(f(x)/g(x)\) 的表达式在 f 与 g 都趋于 0 的那种点处的极限（或“真值”），由这两个函数在所考虑点的邻域内的泰勒展开给出（“洛必达法则”，它很可能应归功于约翰·伯努利）。

撇开这个初等情形不谈，早在 17 世纪末就已经摆在数学家面前的“渐近展开”主要问题，是当 n 很大时形如 \(\sum_{k=1}^{n}f(k)\) 的和的精确或近似计算；这种计算无论对插值、对级数和的数值估计，还是在概率演算中，实际上都是必需的，在概率演算中，诸如 \(n!\) 或 \(\binom{a}{n}\) 这样的“大数的函数”起着举足轻重的作用。早在牛顿，为了在 n 大时得到 \(\sum_{k=1}^{n}\frac{1}{a+k}\) 的近似值，就指出了一种方法，它（在这个特殊情形）归结为计算欧拉–麦克劳林公式的前几项（{[}262{]}，第 II 卷，第 309–310 页）。临近该世纪末，雅各布·伯努利在他关于概率演算的研究中提出要确定和 \(S_{k}(n)=\sum_{p=1}^{n}p^{k}\)，即 n 的多项式\(^{1}\)，他发现了这些多项式的一般构成规律（但未给出证明），从而第一次在这些多项式系数的表达式中引入了以他的名字命名的那些数，以及使这些数得以计算的递推关系（{[}19 b{]}，第 97 页）。1730 年，斯特林通过一种按递推计算系数的程序，对 n 无限增长时的 \(\sum_{k=1}^{n}\log(x+ka)\) 得到了一个渐近展开。

欧拉关于级数及相关问题的决定性工作定年在 1730 至 1745 年之间。设 \(S(n) = \sum_{k=1}^{n} f(k)\)，他把泰勒公式应用于函数 \(S(n)\)，由此得到

\[
f (n) = S (n) - S (n - 1) = \frac {d S}{d n} - \frac {1}{2 !} \frac {d ^ {2} S}{d n ^ {2}} + \frac {1}{3 !} \frac {d ^ {3} S}{d n ^ {3}} - \dots ,
\]

他用待定系数法“反转”这个方程，同时寻找形如下式的解

\[
S (n) = \alpha \int f (n) d n + \beta f (n) + \gamma \frac {d f}{d n} + \delta \frac {d ^ {2} f}{d n ^ {2}} + \dots ;
\]

于是他逐步得到

\[
S (n) = \int f (n) d n + \frac {f (n)}{2} + \frac {1}{12} \frac {d f}{d n} - \frac {1}{720} \frac {d ^ {3} f}{d n ^ {3}} + \frac {1}{30,240} \frac {d ^ {5} f}{d n ^ {5}} - \dots
\]

起初他还不能确定系数的构成规律（{[}108 a{]}，(1)，第 XIV 卷，第 42–72 页及第 108–123 页）。但大约在 1735 年，通过与多项式分解为一次因式的类比，他毫不犹豫地写下了公式

\[
\begin{array}{c} 1 - \frac {\sin s}{\sin \alpha} = \\ \left(1 - \frac {s}{\alpha}\right) \left(1 - \frac {s}{\pi - \alpha}\right) \left(1 - \frac {s}{- \pi - \alpha}\right) \left(1 - \frac {s}{2 \pi + \alpha}\right) \left(1 - \frac {s}{- 2 \pi + \alpha}\right) \dots \end{array}
\]

令整个级数两个成员的展开式的系数相等，他便特别地（对 \(\alpha = \frac{\pi}{2}\)）得到了级数 \(\sum_{n=1}^{\infty} \frac{1}{n^{2k}}\) 的著名表达式\(^2\)，即用 \(\pi\) 的幂表出的那些（同上出处，第 73–86 页）。几年之后，他终于认识到这些 \(\pi\) 的幂的系数与他的求和公式的系数由同样的方程给出，并认出了它们与伯努利所引入的数、与 \(z/(e^z - 1)\) 的级数展开的系数之间的联系（同上出处，第 407–462 页）。

与欧拉相互独立，麦克劳林在大约同一时间、沿一条稍欠冒险而与今天所遵循的路径相近的途径，到达了同一个求和公式：他实际上是在迭代那个把 \(f(x)\) 用差 \(f^{(2k + 1)}(x + 1) - f^{(2k + 1)}(x)\) 表达出来的“泰勒式”公式，而这个公式他是用待定系数法“反转”这些差的泰勒展开而得到的（{[}214{]}，第 II 卷，第 672–675 页）；此外，他也没有注意到为欧拉所发现的系数构成规律。

但是，麦克劳林像欧拉以及他那个时代的所有数学家一样，把他所有的公式都作为级数展开来陈述，其收敛性甚至无人研究。这并不是说收敛级数的概念在当时完全受到忽视：由雅各布·伯努利已经知道调和级数发散，而欧拉本人在用他的求和公式估计这个级数前 n 项之和时也记录了这一结果（{[}108 a{]}，(1)，第 XIV 卷，第 87–100 页及第 108–123 页）；也是欧拉注意到两个相邻的伯努利数之比无限增长，由此可知以这些数为系数的整幂级数不可能收敛（同上出处，第 357 页）。\(^{3}\)但是形式计算的倾向更为强大，欧拉本人那非凡的直觉也不能阻止他有时堕入荒谬，例如当他写下 \(0 = \sum_{n=-\infty}^{+\infty} x^{n}\) 时（同上出处，第 362 页）。\(^{4}\)

我们在别处（见第 153 页）已经说过，19 世纪初的数学家们如何厌倦了这种无节制、无根据的形式主义，而把分析重新引回严格性的道路。一旦收敛级数的概念得到精确化，就出现了对一些简单判别法的需要，用以通过与已知的积分或级数相比较来证明积分和级数的收敛性；柯西在他的《代数分析》中给出了若干判别法（{[}56 a{]}，(2)，第 III 卷），而阿贝尔在一篇身后发表的论文中（{[}1{]}，第 II 卷，第 197–205 页）得到了对数判别法。另一方面，柯西（{[}56 a{]}，(1)，第 VIII 卷，第 18–25 页）澄清了诸如斯特林级数这类级数（由应用欧拉–麦克劳林公式得到，常被称为“半收敛级数”）的悖论：他证明，如果（由于欧拉关于伯努利数的评论）这样一个级数的通项 \(u_k(n)\) 在 \(n\) 取固定值时随 \(k\) 无限增长，那么尽管如此，对固定的 \(k\) 值，部分和 \(s_k(n) = \sum_{h=1}^{k} u_h(n)\) 仍给出级数所“代表”的函数的一个渐近展开（对 \(n\) 趋于 \(+\infty\)），而且 \(k\) 越大越精确。

在古典分析的大部分计算中，都有可能得到函数渐近展开的一般构成规律，其项数可以任意多；这一事实促成了级数与渐近展开之间一种持久的混淆（至少在语言上）；以致 H. 庞加莱在 1886 年费心把渐近展开的初等规则加以编纂时（{[}251 a{]}，第 I 卷，第 290–296 页）（依 \(1/x\) 的整幂、在 \(+\infty\) 的邻域中），使用的仍然是级数论的语言。只是随着来自解析数论的渐近展开的出现，才在渐近展开的概念与级数的概念之间实现了明确的区分，原因在于，在这门理论所处理的大多数问题中，所寻求的展开中只有很少几项（通常只有一项）能够显式地得到。

这些问题也使数学家们熟悉了变量（实变量或整数变量）的幂以外的比较尺度的使用。这一推广应追溯到 P. 杜布瓦–雷蒙的工作{[}94 a and b{]}，他第一个系统地处理了在一点的邻域中比较函数的问题，并在极具独创性的著作中认识到比较尺度的“非阿基米德”特征，同时以一般的方式研究了比较关系的积分与微分，并由此引出大量有趣的结论{[}94 b{]}。然而他的证明缺乏清晰性与严格性，杜布瓦–雷蒙这些结果的正确表述应归功于 G. H. 哈代{[}147{]}：他的主要贡献在于认识并证明了一类“初等函数”即函数 (H) 的存在性，对这类函数，分析的通常运算（尤其是微分）可以应用于比较关系。

